% theme: parallel_lines_and_segment_ratios
% 叩き台。解答生成済み。レビュー用図を挿入済み。
% 表：基本図形を混ぜた小問集合8問。裏：複数の平行線・横切る直線を含む小問集合8問。
\MajorQuestionLesson{平行線と線分の比}
\SetRowSpace{3mm}
\Needspace{90mm}
\TypeQuestion{図で、同じ印をつけた直線は平行である。図中に示された文字の値を求めなさい。長さの単位は $\mathrm{cm}$ とする。}{parallel_segment_ratio_basic_mix}
\QQFigRow[46mm]{fig/lesson_08_basic_04}{$\Lseg{AB}=6,\ \Lseg{BD}=3,\ \Lseg{AC}=8$}{$4$}{fig/lesson_08_basic_02}{$\Lseg{OA}=4,\ \Lseg{OC}=10,\ \Lseg{OB}=6$}{$\dfrac{20}{3}$}
\QQFigRow[46mm]{fig/lesson_08_complex_01}{$\Lseg{AC}=18,\ \Lseg{AB}=6,\ \Lseg{DG}=15$}{$10$}{fig/lesson_08_basic_01}{$\Lseg{AD}=4,\ \Lseg{DB}=6,\ \Lseg{AE}=3,\ \Lseg{BC}=10$}{$x=\dfrac{9}{2},\ y=4$}
\QQFigRow[46mm]{fig/lesson_08_complex_04}{$\Lseg{OB}=5,\ \Lseg{OC}=10,\ \Lseg{CD}=14$}{$7$}{fig/lesson_08_basic_03}{$\Lseg{AB}=4,\ \Lseg{BC}=7,\ \Lseg{DE}=6$}{$\dfrac{21}{2}$}
\QQFigRow[46mm]{fig/lesson_08_basic_07}{$\Lseg{AD}=3,\ \Lseg{AB}=8,\ \Lseg{AE}=4.5,\ \Lseg{BF}=2,\ \Lseg{BC}=12$}{$x=12,\ y=9$}{fig/lesson_08_basic_08}{$\Lseg{AB}=3,\ \Lseg{BC}=5,\ \Lseg{CD}=4,\ \Lseg{EF}=4.5$}{$6$}

\LessonQuestionBreak
\Needspace{90mm}
\TypeQuestion{図で、同じ印をつけた直線は平行である。図中に示された文字の値を求めなさい。長さの単位は $\mathrm{cm}$ とする。}{parallel_segment_ratio_complex_mix}
\QQFigRow[46mm]{fig/lesson_08_basic_05}{$\Lseg{AB}=6,\ \Lseg{CD}=12$}{$8$}{fig/lesson_08_complex_02}{$\Lseg{AB}=4,\ \Lseg{BC}=6,\ \Lseg{CD}=8,\ \Lseg{EH}=27$}{$9$}
\QQFigRow[46mm]{fig/lesson_08_composite_01}{$\Lseg{BC}=6,\ \Lseg{CD}=9,\ \Lseg{EF}=8,\ \Lseg{FP}=4,\ \Lseg{HP}=7.5$}{$x=12,\ y=10$}{fig/lesson_08_complex_05}{$\Lseg{AB}=8,\ \Lseg{CD}=14,\ \Lseg{AE}:\Lseg{ED}=3:2$}{$x=\dfrac{16}{5},\ y=\dfrac{42}{5}$}
\QQFigRow[46mm]{fig/lesson_08_basic_06}{$\Lseg{AD}=6,\ \Lseg{DB}=3,\ \Lseg{OC}=9$}{$6$}{fig/lesson_08_complex_08}{$\Lseg{AE}=4,\ \Lseg{EC}=6,\ \Lseg{BF}=7.5$}{$5$}
\QQFigRow[46mm]{fig/lesson_08_composite_02}{$\Lseg{PF}=6,\ \Lseg{PG}=9,\ \Lseg{PH}=12,\ \Lseg{AB}=10$}{$x=8,\ y=15$}{fig/lesson_08_complex_06}{$\Lseg{AD}=4,\ \Lseg{DF}=3,\ \Lseg{FB}=5,\ \Lseg{BC}=18$}{$\dfrac{21}{2}$}
